% SKELETON ONLY -- NOT \input BY manuscript/BGS_unified.tex AND NOT PART OF THE BUILD.
% Prepared 1 November 2026 as the outline for a future Lines I--N chapter.
% The Line N theorem is expected to move from v4 to v5 after the H1 and H4 research lines
% report; the chapter text will be written only after that. Titles and states below were
% copied from record/claims.json when this file was written; the registry is authoritative.
% Status when written: BGS-TRACE-005 is open; R' (BGS-TRACE-206) is open and not refuted;
% Theorem R#-mod v4 (BGS-TRACE-299) is conditional on the open inputs BGS-TRACE-300..306.
\section{Lines I--N: the polymer-gas route to frozen \texorpdfstring{$R'$}{R'} (skeleton)}
\label{sec:lines-I-N-skeleton}
\emph{Placeholder outline. No statement in this file is asserted beyond its registered state.}

% 114 claims, BGS-TRACE-211 .. BGS-TRACE-324: conditional 3, open 18, proved 89, verification-pending 4
\begin{itemize}
  \item \texttt{BGS-TRACE-211} [proved] Line I Prop I1: relative phase split Z\_m(nu0+delta) = Z\_m(nu0) sigmahat(delta) (k1 kick, any unit nu0)
  \item \texttt{BGS-TRACE-212} [proved] Line I Prop I2/I2': dilation-coset L\textasciicircum{}2 law and per-member Gauss-sum formula for the members Z\_m(nu0 + 2hj/d), d odd, d | h\_m
  \item \texttt{BGS-TRACE-213} [proved] Line I Thm I4: exact reduction of Conjecture R' to (E''), (T'') and the dilation remainder (D)
  \item \texttt{BGS-TRACE-214} [proved] Line I Lemma I5': Fourier expansion of weighted horocycle sums of |Z\_nu|\textasciicircum{}2 (all residues and units via Ramanujan sums)
  \item \texttt{BGS-TRACE-215} [proved] Line I Thm I6 v2: certified bound on the d = 7 dilation square term for the k1 kick at eta = 2, rho = (+-2,8), all m == 3 (mod 6)
  \item \texttt{BGS-TRACE-216} [proved] Line I: exact monomial intertwiners between the eight (m mod 24, a = +-2) cases of the (+-2,8) twisted G2 operators and the beta -> -beta reflection
  \item \texttt{BGS-TRACE-217} [proved] Line I Thm I9: certified top twisted spectrum for d = 7 (classes 1:3, 2:1): three simple eigenvalues with 1.262 < |mu| < 1.45, so r\_tw > sqrt(lambda)/mu0
  \item \texttt{BGS-TRACE-218} [proved] Line I: aliasing (Moebius) identity for the unit-residue mean F\_m
  \item \texttt{BGS-TRACE-219} [proved] Line I: unbounded prime support of h\_m (a new odd prime for every m >= 16; every finite prime set fails for all large m)
  \item \texttt{BGS-TRACE-220} [open] Line I: Conjecture R'\_A (amended mean law with actual dilation-member values) for the k1 kick at eta = 2
  \item \texttt{BGS-TRACE-221} [proved] Line J Prop J-C: exact sector law, exact leading d = 7 coefficient B\_C, exact cancellation of sectors 2, 6 and survival of sectors 0, 4 for rho = (+-2,8)
  \item \texttt{BGS-TRACE-222} [proved] Line J: sector-level antiunitary map sigma(V\textasciicircum{}(-2)\_\{tau,a+4\}) = conj sigma(V\textasciicircum{}(+2)\_\{tau,a\}) (and likewise for E) for the (+-2,8) twisted sector operators, for every twist tau in R/Z and for every non-autonomous product (scope widened at the Lines N2-N4 integration)
  \item \texttt{BGS-TRACE-223} [proved] Line J: certified top-orbit and rest certificates for the d = 7 sector operators: |mu\_2:1| in [1.28359503, 1.28359523], |xi| = |mu* mu\_2:1| >= 2.4586893 > sqrt(lambda)
  \item \texttt{BGS-TRACE-224} [proved] Line J window lemma: for theta in [theta\_lo, theta\_up] subset (0, pi], among N = ceil(2 pi/theta\_lo) consecutive k some cos(alpha + k theta) >= cos(theta\_up/2) and some <= -cos(theta\_up/2)
  \item \texttt{BGS-TRACE-225} [proved] Line J Theorem J: the frozen d = 7, (+-2,8) contribution oscillates with growth >= 1.0184224\textasciicircum{}m: both signs in every 6 consecutive members of each class mod 168 (m == 3 mod 6)
  \item \texttt{BGS-TRACE-226} [proved] Line J3 Lemma L16: lifted trace formula for rho = (+-2,16) restricted to the exact 2-dimensional monomial blocks
  \item \texttt{BGS-TRACE-227} [proved] Line J3: the d = 7 contribution of rho = (+-2,16) (with inherited copies) is <= 1\% of (K\_min/2) xi\_lo\textasciicircum{}m for m >= 243
  \item \texttt{BGS-TRACE-228} [proved] Line J3 inheritance theorem: W\_\{ce, e beta'\} = (n\_e/e\textasciicircum{}2) W\_\{c, beta'\} for coprime odd c, e with ce | h\_m (composite levels carry positive multiples of lower-level class sums)
  \item \texttt{BGS-TRACE-229} [proved] Line J3: prime-level complete-sum bound: |W\_\{q,beta\}| <= 2/q\textasciicircum{}2 for non-isotropic beta, (q-2)/q\textasciicircum{}2 for the 2(q-1) isotropic beta (q || h\_m)
  \item \texttt{BGS-TRACE-230} [proved] Line J3 r(q) lemma: for odd m and an odd prime q, q | h\_m iff r(q) := ord\_q(lambda) divides m
  \item \texttt{BGS-TRACE-231} [conditional] Line J3 Theorem J3 (conditional): frozen R' is false along m = 3p under H-natural and H\_arg; unconditionally, if R' holds the conductor-not-in-\{1,7\} parts must cancel the d = 7 term exponentially on S+
  \item \texttt{BGS-TRACE-232} [open] Line J3 hypothesis H-natural: the conductor-not-in-\{1,7\} part O-natural\_m does not cancel the d = 7 term on infinitely many m in S+
  \item \texttt{BGS-TRACE-233} [open] Line J3 hypothesis H\_arg: arg(xi\_2:1)/2 pi is irrational
  \item \texttt{BGS-TRACE-234} [proved] Line K Lemma K3 (gap lemma): for b' >= 447, b' not in \{504, 560, 616, 840\}, every unit class mod gcd(56, b') has maximal cyclic gap < A b' among the units mod b' in that class (A = arccos(0.84431)/pi)
  \item \texttt{BGS-TRACE-235} [proved] Line K Theorem K3: S+ contains infinitely many m = 3p with p == p0 (mod 56) for every unit p0 mod 56, unconditionally (H\_arg removed)
  \item \texttt{BGS-TRACE-236} [proved] Line K Corollary K3: frozen R' implies phi\textasciicircum{}-1 O-natural\_m <= -0.99 K\_j 1.0184224\textasciicircum{}m + 4m + o(m) on an infinite set of m = 3p (each unit class p0 mod 56); hence H-natural alone refutes frozen R'
  \item \texttt{BGS-TRACE-237} [proved] Line K Theorem K1 (light-cone identity): for a prime q != 7, q || h\_m, 2 a square mod q, sum\_\{beta != 0\} Zhat\_beta W\_\{q,beta\} = D\_iso - E\_b/q
  \item \texttt{BGS-TRACE-238} [proved] Line K Prop K2 (large levels by positivity): for m = 3p >= 495 the levels d >= L\_m = (mu*\textasciicircum{}2/xi)\textasciicircum{}\{m(1+eta\_m)\} contribute >= -0.01 (K\_min/2) xi\_lo\textasciicircum{}m; S\_Z <= 0.5000001; eta\_m <= (2.010 + o(1))/log(2p+1)
  \item \texttt{BGS-TRACE-239} [open] Line K hypothesis H-natural\_mid: for some delta > 0 and infinitely many m in S+, O-natural\_m\textasciicircum{}\{<L\} >= -(1 - delta) 0.98 (K\_j/2) xi\_lo\textasciicircum{}m
  \item \texttt{BGS-TRACE-240} [conditional] Line K Theorem K (conditional): frozen R' is false under H-natural or under the alternative hypothesis H-natural\_mid
  \item \texttt{BGS-TRACE-241} [proved] Line K Lemma W1 (Weil bound over the orbit subgroup): for q prime, H = <l> subset F\_q\textasciicircum{}* of order r and A, B in F\_q not both 0, |sum\_\{h in H\} e((Ah + B/h)/q)| <= 2 sqrt q (<= sqrt q if AB = 0); degeneracy criterion
  \item \texttt{BGS-TRACE-242} [proved] Line K Lemma W2' (non-cyclic block AM-GM, every bounded read-out): |Tr(J X)| <= C\_J max(1, sup||L||)\textasciicircum{}\{2k\} e\textasciicircum{}2 max(1, (1/m) sum\_\{i<m\} ||P(i)||\_HS\textasciicircum{}2)\textasciicircum{}\{(Q-1)/2\}
  \item \texttt{BGS-TRACE-243} [proved] Line K Lemma W3 (orbit average): (1/r) sum\_\{i<r\} f\_k(beta lambda\textasciicircum{}i/q) <= I\_k + 2 sqrt q A\_k/r + T\_k(q) for q || h, beta != 0, with T\_k(q) -> 0 (6 sqrt q for conductor 7q)
  \item \texttt{BGS-TRACE-244} [proved] Line K operator input (OP\_k, delta): the cat-map second moments satisfy (I\_k\textasciicircum{}F)\textasciicircum{}\{1/2k\} <= (1 - delta) 1.28359 for every (+-2,8) family and <= (1 - delta) 1.4742 for every (+-2,16) family
  \item \texttt{BGS-TRACE-245} [proved] Line K Theorem W (OP proved in Line L; under the standing inputs): for m = 3p in S+ with m >= m1 (effective, not computed), the conductors q, 7q with q prime, q | h\_\{3p\}, q <= c\_W p\textasciicircum{}2 (c\_W = 10\textasciicircum{}-25.62, q0 = 10\textasciicircum{}13.2) contribute o(xi\_lo\textasciicircum{}m) to O-natural, including inherited copies
  \item \texttt{BGS-TRACE-246} [proved] Line K Theorem K2' (OP proved in Line L; under the standing inputs): for m in S+ with m >= m1 (effective, not computed), frozen R' implies that some open conductor c of h\_m carries a slow lambda-orbit
  \item \texttt{BGS-TRACE-247} [open] Line K Conjecture C1: for the true twisted operator families, lim\_k (I\_k\textasciicircum{}F)\textasciicircum{}\{1/k\} = 1
  \item \texttt{BGS-TRACE-248} [proved] Line L collision structure: alpha\_v = sum\_j v\_j lambda\textasciicircum{}j = 0 iff (X\textasciicircum{}2 - 6X + 1) | v(X); carries and the Green's-function bound
  \item \texttt{BGS-TRACE-249} [proved] Line L Lemma O1 (word Hilbert-Schmidt bound) for the beta-weighted lattice-lifted k1 families, with the outermost-letter tau assignment
  \item \texttt{BGS-TRACE-250} [proved] Line L Theorem OP\_abs (cancellation-free second moment): I\_k <= C\_pair S\_k and A\_k <= C\_pair prod\_j A(tau\_j)\textasciicircum{}2 for the true lattice-lifted k1 families
  \item \texttt{BGS-TRACE-251} [proved] Line L OP certificate: certified I\_k\textasciicircum{}\{1/2k\} <= 1.26278 / 1.13642 / 1.39243 / 1.30276 for (+-2,8)V/E and (+-2,16)V/E, below the targets 1.28359 and 1.4742
  \item \texttt{BGS-TRACE-252} [proved] Line L T\_k bound and Weil-range constants: c\_W = 10\textasciicircum{}-25.62 <= min\_F c\_F, q0 = 10\textasciicircum{}13.2; the range q <= c\_W p\textasciicircum{}2 is non-vacuous only from p \textasciitilde{} 8.3e25
  \item \texttt{BGS-TRACE-253} [proved] Line L (+-2,16) narrowing: the (+-2,16) levels are absorbed for q >= (1.6677\textasciicircum{}2/xi\_lo)\textasciicircum{}\{m(1+eta)\}, eta <= (4.282 + o(1))/log(2p+1), so the (+-2,16) open range is c\_W p\textasciicircum{}2 < q < (1.6677\textasciicircum{}2/xi\_lo)\textasciicircum{}\{m(1+eta)\} (1.6677\textasciicircum{}2/xi\_lo = 1.13118...)
  \item \texttt{BGS-TRACE-254} [proved] Line L Theorem U (level-h scalar L\textasciicircum{}2 unitarity): the TRACE-206 autocorrelation R(v) vanishes at every exact collision and every collision mod h, so sum\_x |A(x)|\textasciicircum{}2 = 1 and sum\_\{beta mod h\} |Zhat(beta/h)|\textasciicircum{}2 = h\textasciicircum{}2
  \item \texttt{BGS-TRACE-255} [proved] Line L2 decomposition: D\_q = (1/q) sum\_\{beta in l(q) minus 0\} Zhat\_beta, D\_iso = D\_q + D\_qbar - (1/q\textasciicircum{}2) sum\_\{beta != 0\} Zhat\_beta (open prime levels q || h\_m)
  \item \texttt{BGS-TRACE-256} [open] Line L2 hypothesis H\_mean: for some delta > 0 and infinitely many m in S+, |D\_iso(q) - E\_b(q)/q| <= ((1 - delta) target)\textasciicircum{}m at every open prime level q, with the analogue at open composite levels
  \item \texttt{BGS-TRACE-257} [proved] Line L2: H\_mean implies that frozen R' is false along m = 3p (a proved implication; H\_mean is open)
  \item \texttt{BGS-TRACE-258} [proved] Line L2 Prop S: beta = 1 is slow at every prime q || h\_m with log\_lambda q >= C\_k + 6.1e-6 m (C\_k effective, not computed); consequence: K2' gives no contradiction at m where h\_m has such a prime below L\_m
  \item \texttt{BGS-TRACE-259} [proved] Line L2 Lemma A1 (exact second moment mod q): sum\_\{beta in l(q)\} |Zhat\_beta|\textasciicircum{}2 = q sum\_u |A(u + q)|\textasciicircum{}2, |D\_q|\textasciicircum{}2 <= (1 - 1/q) C\_q, C\_q = 1 + sum\_\{alpha\_v in q minus 0\} R(v) - |Z|\textasciicircum{}2/q
  \item \texttt{BGS-TRACE-260} [proved] Line L2 Prop A2 (phase-blind barrier): every cancellation-free collision bound for C\_q + |Z|\textasciicircum{}2/q is >= G\textasciicircum{}\{2m\}/q, G = sum\_n |J\_n(1.8)| = 2.3693986; it certifies the target only for q >= 3.4074\textasciicircum{}m ((+-2,8)) or 2.5832\textasciicircum{}m ((+-2,16)), above h\_m itself
  \item \texttt{BGS-TRACE-261} [proved] Line L2 Lemmas B1-B3 (height/norm, split q): q <= 1.4571 lambda\textasciicircum{}\{r-1\} ||F||\_inf\textasciicircum{}2 for sporadic folded collisions; Minkowski vectors with ||F||\_inf <= q\textasciicircum{}\{1/r\}; the sum over prime ideals above h
  \item \texttt{BGS-TRACE-262} [proved] Line L2: for all sufficiently large p the height zone of Lemma B1 and the open range (c\_W p\textasciicircum{}2, L\_m) are disjoint, for (+-2,8) and (+-2,16)
  \item \texttt{BGS-TRACE-263} [proved] Line L3 S1 (walk-measure form): A = chi B\_m with chi(x) = e(-nu0 y' Nm x/2h) and Bhat\_m(y) = prod\_\{j<m\} g(Tr lambda\textasciicircum{}j y), g(t) = e\textasciicircum{}\{i 1.8 cos 2 pi t\}, so |Bhat\_m| = 1
  \item \texttt{BGS-TRACE-264} [proved] Line L3 S2 (CRT split): A(q) is a level-2h/q Gauss sum of B\_m on q; |A(q)| <= h/q and |D\_q| <= 2h/q when |Z| <= h, which gives |D\_q| <= T\textasciicircum{}m only for q >\textasciitilde{} 1.8808\textasciicircum{}m, above L\_m
  \item \texttt{BGS-TRACE-265} [proved] Line L3 S3 (no Gauss-side folding): h\_\{3p\} = h\_p (4 h\_p\textasciicircum{}2 + 3); every prime of K\_p = 4 h\_p\textasciicircum{}2 + 3 other than 7 (rank 3) has rank 3p, so the level-K\_p factor of the m = 3p phase does not reduce to m = p
  \item \texttt{BGS-TRACE-266} [proved] Line L3 S4 (transfer-side spectral form at rank-p levels): if r(c) = ord\_c(lambda) = p for an odd conductor c (prime or composite, e.g. c = h\_p) then Zhat\textasciicircum{}\{(3p)\}\_beta = c\_\{3p\} Tr\_K[(I\_K tensor <1|) Pi\_beta\textasciicircum{}3 (I\_K tensor |0>)], Pi\_beta the p-step twisted transfer operator (given the Line G/I transfer law)
  \item \texttt{BGS-TRACE-267} [proved] Line L3 Lemma Q (exact line quadrature): (1/q) sum\_\{beta in l(q)\} F(beta/q) = sum\_\{gamma in q\} c\_gamma; so the isotropic-line mean block energy lies in [I\_k - T\_k(q), I\_k + T\_k(q)] for every q, without Weil
  \item \texttt{BGS-TRACE-268} [proved] Line L3 Lemma E (cusp excursions): if dist(Tr(beta lambda\textasciicircum{}j)/q, Z) <= eps for k >= 3 consecutive j (0 < eps <= 1/16, beta != 0), then lambda\textasciicircum{}\{k-1\} <= 264 eps\textasciicircum{}2 q (beta isotropic) and <= 264 eps\textasciicircum{}2 q\textasciicircum{}2 (every beta != 0)
  \item \texttt{BGS-TRACE-269} [proved] Line M M1: gcd(h\_m, h\_n) = h\_gcd(m,n) for odd m, n; for m = 3p every prime q != 7 dividing h\_m has r(q) in \{p, 3p\}, so the open conductors at distinct m = 3p are pairwise disjoint (only 7 is shared)
  \item \texttt{BGS-TRACE-270} [proved] Line M M2 (averaging lemma): a positive-weight averaged form of H-natural (with delta) implies H-natural with delta/2; averaged hypotheses are never weaker
  \item \texttt{BGS-TRACE-271} [proved] Line M M3 (inherited-level no-go for positivity methods): for m = 3p >= 495 and rho = (+-2,8), the level d = h\_p is included and open, and the Prop K2 per-level floor there is >= c 1.13287\textasciicircum{}m, exceeding the D7 demand by more than 1.11238\textasciicircum{}m
  \item \texttt{BGS-TRACE-272} [proved] Line M M4a: the four P* families are disjoint for every odd m (no unit lies in two families at levels <= sqrt h, primaries included; each family has exactly one primary)
  \item \texttt{BGS-TRACE-273} [proved] Line M M4b: for odd m >= 15 the (+-2,8) inclusion rule of the frozen definition includes every level d | h with 1 < d <= sqrt h (no borderline level), on the Line F certified bound
  \item \texttt{BGS-TRACE-274} [open] Line M M4c (open): the (+-2,16) inclusion rule degenerates (no borderline level) for m >\textasciitilde{} 85
  \item \texttt{BGS-TRACE-275} [proved] Line M M5 (refinement of Prop K2): at an included level d, squares offset a cross deficit X only if d >\textasciitilde{} 4|Z\_rho|\textasciicircum{}2/X, and minimising over the member sum gives exactly Prop K2's per-level bound -n\_d |Z\_rho|\textasciicircum{}2/d\textasciicircum{}2 (sharp within the quadratic relaxation)
  \item \texttt{BGS-TRACE-276} [open] Line M M6 (open, conditional on C1-C4): the +-rho pairing law at all levels (cross\_+ + cross\_- = 2 Re sum\_\{a,b\} conj(c\_a z\_a) c\_b y\_b (1 + s\_a s\_b))
  \item \texttt{BGS-TRACE-277} [proved] Line M M7a (arithmetic): ord\_c(lambda) | 3p for c | h\_\{3p\}; lambda\textasciicircum{}p = 1 + 2 h\_p eps\textasciicircum{}p, so lambda\textasciicircum{}p == 1 mod 2h\_p; p = 2 log\_lambda h\_p + 2 log\_lambda 2 + o(1); open prime conductors q | h\_\{3p\} have ord\_q(lambda) <= 3p < q
  \item \texttt{BGS-TRACE-278} [proved] Line N N0 (reduction and budget): the frozen orbit average at a prime conductor q | h\_\{3p\} is I\_k + sum\_\{gamma != 0\} c\_gamma S(gamma beta) with S a binomial sum over H = <lambda mod q> of order r in \{p, 3p\}; no Gauss-type phase survives; per-level threshold eps\_F = eta\_F/(3 A\_k\textasciicircum{}F) independent of m (eps = 10\textasciicircum{}-12.51 for (+-2,8)V)
  \item \texttt{BGS-TRACE-279} [proved] Line N N1.1-N1.3: moment inequality |S|\textasciicircum{}\{2kl\} <= q\textasciicircum{}2 tau\textasciicircum{}\{2kl-2k-2l\} J\_k J\_l, energies J\_2 <= 2 tau\textasciicircum{}2 - tau and J\_3 <= 17.31 tau\textasciicircum{}\{11/3\} (q >= c\_W p\textasciicircum{}2, tau >= 10\textasciicircum{}45), and the explicit binomial bounds (2,2), (2,3), (3,3) with the three-branch isotropic bound (Gauss period and both Cochrane-Pinner branches)
  \item \texttt{BGS-TRACE-280} [proved] Line N Theorem W+ (explicit beyond-Weil range): Theorem W holds, under its standing inputs and for m >= m1 (effective, not computed), for prime conductors q and 7q with q <= max(c\_W p\textasciicircum{}2, c\_N p\textasciicircum{}\{7/3\}), c\_N = 10\textasciicircum{}-113.80; beyond Weil only for p >= 10\textasciicircum{}264.5
  \item \texttt{BGS-TRACE-281} [proved] Line N N2 (qualitative extension, prime conductors): for every fixed C, Theorem W holds for prime conductors q, 7q <= p\textasciicircum{}C with p >= p0(C) = eps\_F\textasciicircum{}\{-1/delta(C)\}, p0(C) effective, constants not computed (Bourgain, JAMS 2005)
  \item \texttt{BGS-TRACE-282} [open] Line N (open [H]): Theorem W for composite conductors c | h\_\{3p\} with c <= p\textasciicircum{}C (at least two prime factors >= 2p + 1) via Bourgain's finite-ring theorem
  \item \texttt{BGS-TRACE-283} [proved] Line N Theorem N3 (scarred conductors carry slow beta): for F = (+-2,8)V, theta\_F = 5.39e-6 and t\_F effective, not computed, every conductor c with r = ord\_c lambda > 2(N + k) and log\_lambda c >= theta\_F (1 + kappa) r/2 + t\_F + k has the slow beta = 1; in particular c = h\_p for every p >= p0 (effective, not computed)
  \item \texttt{BGS-TRACE-284} [proved] Line N Corollary N3': for every m = 3p in S+ with p >= p0 (effective, not computed) the conclusion of K2' holds without assuming frozen R' (c = h\_p is open for the (+-2,8) families and carries the slow beta = 1 in (+-2,8)V), so K2' is never a refutation route there and the unweighted 'exclude slow lambda-orbits' target is false
  \item \texttt{BGS-TRACE-285} [proved] Line N N4 (identification of the inherited level): the conductor-h\_p points ((1/h\_p)O + O\textasciicircum{}dual)/O\textasciicircum{}dual (h\_p\textasciicircum{}2 points) form a subgroup of index 4 in Fix(T\textasciicircum{}p) (4 h\_p\textasciicircum{}2 points), so level-h\_p twisted sums are periodic-orbit sums of the cat map
  \item \texttt{BGS-TRACE-286} [proved] Line N2 Theorem A (beta-weighted K2'): the Omega-hypothesis (BGS-TRACE-287) implies H-natural\_mid and hence that frozen R' is false along m = 3p (a proved implication; the Omega-hypothesis is open)
  \item \texttt{BGS-TRACE-287} [open] Line N2 Omega-hypothesis (open): for infinitely many m in S+ with m >= m1 and every family F, the beta-weighted pressure sum Omega\textasciicircum{}F\_m(k,a) is <= e\textasciicircum{}\{o(m)\} ((1 - delta) T\_F)\textasciicircum{}m for some fixed k, a, delta > 0
  \item \texttt{BGS-TRACE-288} [proved] Line N2 Theorem B: at c = h\_p the class sums are periodic-orbit sums: sum\_beta W\_\{h\_p,beta\} Zhat\_beta = sum\_\{y in G\_p\} W(y) sum\_F a\_F Tr(J\_F M\_F(y)\textasciicircum{}3) over the index-4 subgroup G\_p = (1/h\_p)Z\textasciicircum{}2/Z\textasciicircum{}2 of Fix(A\textasciicircum{}p), with Fourier form and a Bowen-ball comparison with Lebesgue
  \item \texttt{BGS-TRACE-289} [proved] Line N2 Theorem C (pressure reduction): limsup\_p (1/p) log E\_\{y in G\_p\} |Zhat\textasciicircum{}F(y)| <= log Lambda\textasciicircum{}F\_\{k'\} + log r\_2\textasciicircum{}F(k) + eta, with Lambda the k'-step cocycle norm in the O1 space (or in the Line F space via Lemma I) and r\_2 a Collatz-Wielandt bound for the carry automaton
  \item \texttt{BGS-TRACE-290} [proved] Line N2 Proposition E (finite tail-free collision majorant) with arb certificates: r\_2(4000) <= 1.0529012 (8V), 1.0519234 (8E), 1.0534557 (16V), 1.0529202 (16E)
  \item \texttt{BGS-TRACE-291} [proved] Line N3 Lemma 1.1 (geometric twist sequences): for prime q = qq' split, q | h\_p, the isotropic beta in the ideal q have tau\_i(beta) = s l\textasciicircum{}i / q mod 1 with l = lambda mod qbar of order exactly p; unit beta have tau\_i = (s1 l\textasciicircum{}i + s2 l\textasciicircum{}-i)/q
  \item \texttt{BGS-TRACE-292} [proved] Line N3 Lemma 1.2 (isotropic part as a line quadrature): at a prime q || h, q | h\_p, the isotropic cross part is 2 Re(conj Z\_rho ((q - 2)/q)(D\_q + D\_qbar)), the same weight for +-rho
  \item \texttt{BGS-TRACE-293} [proved] Line N3 M6 condition C3: W\_\{q,beta\}(-rho) = conj W\_\{q,beta\}(rho) at every prime level q || h\_m for rho = (+-2,8), and at composite levels via inheritance (isotropic weights equal)
  \item \texttt{BGS-TRACE-294} [proved] Line N3 Lemma I / Line N4 Lemma I16 (intertwining): sup\_y ||P\textasciicircum{}\{O1\}\_\{k'\}(y)||\textasciicircum{}\{1/k'\} <= (C\_I kappa\_F)\textasciicircum{}\{1/k'\} Lambda\_free\textasciicircum{}F with C\_I independent of k', p and the twists, for all four families
  \item \texttt{BGS-TRACE-295} [proved] Line N4 certified free-family joint spectral radii [CA]: Lambda\_free\textasciicircum{}\{8V\} <= 1.9656353, Lambda\_free\textasciicircum{}\{8E\} <= 0.4915932, Lambda\_free\textasciicircum{}\{16V\} <= 2.9413897, Lambda\_free\textasciicircum{}\{16E\} <= 0.4474396, uniform over all twist sequences
  \item \texttt{BGS-TRACE-296} [proved] Line N3/N4 Theorem C': for p with h\_p squarefree and p >= p0 (effective, not computed), the unit-beta cross part at c = h\_p is O(((1 - delta) T\_F)\textasciicircum{}m) for all four families, with unit-point rates (Lambda r\_2)\textasciicircum{}\{1/3\} <= 1.27438 (8V), 0.80266 (8E), 1.45789 (16V), 0.77812 (16E)
  \item \texttt{BGS-TRACE-297} [proved] Line N4 gauge lemma: L\_tau T\_sigma = T\_\{lambda sigma\} L\_\{tau - kappa lambda sigma\}, L\_tau G\_nu = G\_\{nu/lambda\} L\_\{tau + nu/lambda\}; a product along a two-sided geometric twist t0 + a lambda\textasciicircum{}i + b lambda\textasciicircum{}-i is a gauge conjugate of L\_\{t0\}\textasciicircum{}n up to a unimodular scalar
  \item \texttt{BGS-TRACE-298} [open] Line N3/N4 open item: the isotropic part at q | h\_p (and at the mixed levels 7q), i.e. a bound |D\_q + D\_qbar| <= ((1 - delta) T)\textasciicircum{}m for the isotropic lines at every prime q | h\_p, together with composite d | h\_p having an isotropic prime component
  \item \texttt{BGS-TRACE-299} [conditional] Line N Theorem R\#-mod v4 (= v3 after lineN16 edits E1-E7) (conditional): if H0-H8 (BGS-TRACE-300..306 and the standing package) hold at p and B\_p >= kappa > 0 for infinitely many primes p >= p\_thr, then W'\_\{3p\}/m -> infinity and frozen R' is false along m = 3p (conditional; not a refutation)
  \item \texttt{BGS-TRACE-300} [open] Line N hypothesis H1 (open [H]): collision lemma - lattice-to-collision reduction for shell sums with constant C <= 0.082 at J0 = 12, plus the A\_bad tail j > 12
  \item \texttt{BGS-TRACE-301} [open] Line N hypothesis H3 (open): Kotecky-Preiss polymer weight e\textasciicircum{}\{a|gamma|\} with a = 5.93e-6 at g0 = 4 (log -> modulus), finite part certified, tail j > 12 = H1
  \item \texttt{BGS-TRACE-302} [open] Line N hypothesis H4 (open, no working route): core -> true conversion of the interaction constant, now = difference-form (pair-minus-product) conversion + the G > 15 true-functional tail
  \item \texttt{BGS-TRACE-303} [open] Line N hypothesis H5 (open): composite h\_p - (a) mixed classes (R-mix) E\_\{mixed,d\} F\_p <= e\textasciicircum{}\{o(p)\} d\textasciicircum{}c r\_2\textasciicircum{}p with c < 0.387 uniformly in d | h\_p; (b) the 2\textasciicircum{}omega - 2 cross lines; (c) partially isotropic points; or p in P = \{h\_p prime\}
  \item \texttt{BGS-TRACE-304} [open] Line N hypothesis H6 (open): inputs of Theorem C' - an explicit p0 and squarefreeness of h\_p (or conductor-q\textasciicircum{}a weights)
  \item \texttt{BGS-TRACE-305} [open] Line N hypothesis H7* (open): for included levels d not in \{h\_p\} u \{e h\_p\}, d < L\_fl = h\_p\textasciicircum{}\{1+o(1)\}: |Lin\_d| <= eps\_d |Z\_rho|\textasciicircum{}2/h\_p with sum eps\_d <= kappa0, i.e. cross-correlation growth < T' = 1.42785; needed even for prime h\_p
  \item \texttt{BGS-TRACE-306} [open] Line N hypothesis H8 (open): deep wrap-around points of gap (a) at composite h\_p - free if per-point weights are not d-amplified, otherwise need c < 1.1718 l0 / ln d (= 0.6647 l0 / log\_lambda d)
  \item \texttt{BGS-TRACE-307} [proved] Line N H2 met [CA-fp]: core interaction constant delta\_int\textasciicircum{}core <= 5.14476e-7 (c3\_mid, closed F14.1; referee arb 5.1447502e-7)
  \item \texttt{BGS-TRACE-308} [proved] Line N [CA-fp]: rhobar <= 0.28 by an eigen-solver-free residual certificate (N15 b\_cp3)
  \item \texttt{BGS-TRACE-309} [proved] Line N11 coding lemma [P]: the isotropic sum at q = h\_p is coded by the gas, X\_coded = 2[(1 + Lambda)\textasciicircum{}p e\textasciicircum{}\{p Delta\} - 1] - (dense part), and |1 + X\_gas| >= 2(1 + Re Lambda)\textasciicircum{}p e\textasciicircum{}\{-p delta\} - 1 given the model form
  \item \texttt{BGS-TRACE-310} [proved] Line N [CA]: Re Lambda\_lo >= 6.238e-4 (Re Lambda\_lo' = 6.234e-4 after charges)
  \item \texttt{BGS-TRACE-311} [proved] Line N13 adapted norm [CA] (cond\_H 21.103, Lip\_H 167.213, Qtilde 0.45) and quadratic step bound N\_H <= 1 + 1.01e-12 + 4260 t\textasciicircum{}2 for |t| <= 1/256 [CA]+[P]
  \item \texttt{BGS-TRACE-312} [proved] Line N15 exact carry reduction [P]: |e| <= 15 (sharp |e| <= 14, 29 values, 633 reachable pair states) and the torus-mean identity [P]
  \item \texttt{BGS-TRACE-313} [proved] Line N7/N8b Lemmas P, P' (all isotropic points at q = h\_p, with the 4-fixed-point correction N8b-F3), P'\_c, the classification and the weight amplification A\_c [P]
  \item \texttt{BGS-TRACE-314} [proved] Line N7: the (+-2,16) contribution is closed at all levels for p >= 339 (two-term B16 form; crude form p >= 498) [CA]+[P]
  \item \texttt{BGS-TRACE-315} [proved] Line N17b K\_p arithmetic [P]: h\_\{3p\} = h\_p K\_p with K\_p = 4 h\_p\textasciicircum{}2 + 3, 7 | K\_p, gcd(h\_p, K\_p) | 3; h\_15 = 7 * 31\textasciicircum{}2 * 41
  \item \texttt{BGS-TRACE-316} [proved] Line N17b CRT count [P]: at h = 8119 = 23 * 353 (m = 33) S = CRT and n = 7371 = 21 * 351; 7743 - 7371 = 372 are the non-unit nu excluded by BGS-TRACE-206
  \item \texttt{BGS-TRACE-317} [proved] Line N17b level bookkeeping [P]: W'\_m = NR part + phi\textasciicircum{}-1 sum Dev\_d(rho) (BGS-TRACE-213), NR >= -RMT >= -4m, level floor -(h\_p/d) S\_rho, inclusion m0 = 15
  \item \texttt{BGS-TRACE-318} [proved] Line N18b/N19b gas vs Robin crossover [P-arith], an arithmetic implication: under the idealisation tails = Y = X\_cross = eps = 0 and n/h\_p := 1, the gas lower bound exceeds the Robin bound on inherited levels for all p >= 1386 at the core value delta = 5.14476e-7 (BGS-TRACE-307) and, for the INPUT value delta = 2.470e-6, for all p >= 1390; first prime 1399
  \item \texttt{BGS-TRACE-319} [proved] Line N14b/N18b Lemma 1 (sublattice separation) [P]: for squarefree d with split prime factors, delta-close pairs of the conductor-d components H\_d separate after n\_* steps, with delta0 = 1/(4 sqrt2 (lambda+1)) = 0.025888 and constant C = 16 lambda\textasciicircum{}2 delta\textasciicircum{}2 h\textasciicircum{}2 lambda\textasciicircum{}\{-p\}
  \item \texttt{BGS-TRACE-320} [verification-pending] Line N14b Lemma 2 (subgroup C', c\_B = 0.7375) [P-sketch, asymptotic, constants ignored], covering iso parts d <= 1.588\textasciicircum{}p; p = 31 covered (H5, not P)
  \item \texttt{BGS-TRACE-321} [proved] Line N15 [CA-fp]: A\_bad finite part j <= 12 = 2.031569e-6 at g0 = 4 (c3 radii) and short types sum |w| = 9.970e-7 (173 rows incl. alpha = 4)
  \item \texttt{BGS-TRACE-322} [verification-pending] Line N16 [CA-fp]+[P-sketch]: delta\_true(G <= 15) <= 2.470e-6 (vector-chain conversion per window, E/|r| \textasciitilde{} 1.7e-3)
  \item \texttt{BGS-TRACE-323} [verification-pending] Line N15 [CA-fp]: SVD-free step-norm table for |t| > 1/256, max 2.91656
  \item \texttt{BGS-TRACE-324} [verification-pending] Line N8 Prop X' (explicit error, p >= 1289) with the 15 small-depth classes [CA]: [P mod Line F certificates + class CA] for the certified classes; k-parity mechanism and E-family absorption [P-sketch]
\end{itemize}
